%-------------------------------------------------------------------------------
% CONFIGURATIONS
%-------------------------------------------------------------------------------
\documentclass[11pt, a4paper]{awesome-cv}
\geometry{left=0.8cm, top=1cm, right=0.8cm, bottom=0.5cm, footskip=.5cm}
\fontdir[fonts/]
\colorlet{awesome}{awesome-red}
\setbool{acvSectionColorHighlight}{true}
\renewcommand{\acvHeaderSocialSep}{\quad|\quad}

%-------------------------------------------------------------------------------
% INFORMATIONS PERSONNELLES
%-------------------------------------------------------------------------------
% \photo{img/cv_pc.png}
\name{Yassire}{Ammouri}
\position{Ingénieur IA \space | \space Diplômé en Data Science \& Ingénierie}
\address{Salé - Rabat, Maroc}
\mobile{(+212) 648289422}
\email{yassire.ammouri@um5r.ac.ma}
\github{Yassire1}
\linkedin{yassire-ammouri}

%-------------------------------------------------------------------------------
\begin{document}
\makecvheader

%-------------------------------------------------------------------------------
\vspace{-0.8cm}
\cvsection{Profil}

\begin{cvparagraph}
Étudiant en dernière année de Data Science \& Ingénierie à la recherche d'un premier poste de Data Scientist. Maîtrise de \textbf{Python} (pandas, scikit-learn, TensorFlow), \textbf{SQL} et la visualisation de données. Expérience en pipelines ML de bout en bout, ingénierie des caractéristiques, deep learning pour les données temporelles et suivi d'expériences avec Weights \& Biases. Solide base en statistiques et manipulation de données. Désireux de contribuer à la prise de décision basée sur les données dans un environnement collaboratif.
\end{cvparagraph}

%-------------------------------------------------------------------------------
\vspace{-0.6cm}
\cvsection{Formation}
\vspace{0.2cm}
\cventry
  {SUPMTI Rabat}
  {Diplôme d'Ingénieur en Data Science \& Informatique}
  {Sept. 2023 -- Présent }
  {Rabat, Maroc}
  {}
  
\vspace{-0.5cm}
\cventry
  {Faculté des Sciences, Tétouan}
  {Licence en Électronique }
  {Sept. 2020 -- Août 2023}
  {Tétouan, Maroc}
  {}
  
\vspace{-0.5cm}
\cventry
  {Lycée Mohamed 5}
  {Baccalauréat en Physique}
  {Sept. 2019 -- Août 2020}
  {Ouazzane, Maroc}
  {}
  
\vspace{-0.4cm}
%-------------------------------------------------------------------------%
\cvsection{Expérience Professionnelle}

\vspace{0.1cm}
\cventry
  {OCP SOLUTIONS}
  {Stagiaire Ingénieur IA Backend }
  {Fév 2026 -- Présent}
  {Rabat, Maroc}
  {
    \begin{cvitems}
      \item {Développement d'un agent LangGraph stateful pour la génération de micro-cours pilotée par l'IA avec des états typés, un routage conditionnel, une confirmation avec humain dans la boucle et des boucles de rétroaction adaptatives qui ajustent la difficulté des leçons en fonction des performances de l'utilisateur.}
      \item {Mise en place d'un chat SSE en temps réel via FastAPI, servant des réponses LLM au niveau du token tout en extrayant des métadonnées JSON structurées et en déclenchant la planification asynchrone des programmes en arrière-plan.}
      \item {Développement d'un pipeline TTS multimodal avec DashScope CosyVoice, segmentation markdown, blocs audio synchronisés et génération en arrière-plan par lots pour une livraison de leçons texte-audio synchronisée.}
      \item {Réduction de la latence au démarrage de 20s+ à <2s grâce à l'initialisation anticipée de LangGraph et des connexions LLM, permettant des conversations stateful concurrentes évolutives avec un raisonnement transparent via des événements de streaming.}
      \item {Orchestration de 6 microservices FastAPI avec Docker Compose, PostgreSQL, Redis et un frontend Next.js 16, atteignant 88\% de couverture de tests sur le service du moteur IA.}
    \end{cvitems}
  }

\vspace{-0.3cm}
\cventry
  {Laboratoire LRIT - Faculté des Sciences, Université Mohammed V}
  {Stagiaire Assistant de Recherche -- Data Scientist}
  {Fév 2025 -- Sep 2025}
  {Rabat, Maroc}
  {
    \begin{cvitems}
      \item {Développement d'un \textbf{cadre d'apprentissage fédéré centré sur les données} pour la prédiction de pannes éoliennes multiclasses (6 classes : Normal + 5 pannes de composants), traitant plus de 400K+ enregistrements SCADA multivariés de séries temporelles sur 4 éoliennes sur 2 ans.}
      \item {Conception de pipelines de prétraitement de données de bout en bout incluant l'étiquetage temporel dynamique (fenêtres de 30/60 jours et spécifiques aux composants), la sélection hybride de caractéristiques (120+ caractéristiques brutes $\rightarrow$ 26 prédicteurs métier) et la détection de fenêtres post-panne avec logique de fallback hiérarchique.}
      \item {Mitigation du déséquilibre extrême des classes ($<$2\% d'échantillons de pannes) via les fonctions de pondeur de classes, l'optimisation de la perte focale, le suréchantillonnage synthétique basé sur TimeGAN et le sous-échantillonnage temporel intelligent.}
      \item {Développement d'une architecture deep learning hybride \textbf{1DCNN-BiLSTM}, réalisant plus de 10 expériences comparatives sur les configurations de modèles, les fonctions de perte et les variantes de pipelines de données.}
      \item {Implémentation de l'apprentissage fédéré utilisant le framework \textbf{Flower} avec les stratégies d'agrégation FedAvg/FedProx, déployé via \textbf{Docker Compose} pour simuler un entraînement collaboratif préservant la confidentialité sur 3 clients éoliens distribués.}
      \item {Utilisation de \textbf{Weights \& Biases (W\&B)} pour le suivi des expériences, la journalisation des hyperparamètres et la documentation reproductible de tous les entraînements, métriques et artefacts.}
    \end{cvitems}
  }

%-------------------------------------------------------------------------------
\vspace{-0.4cm}
\cvsection{Compétences Techniques}
\vspace{0.1cm}
\begin{cvskills}
    \cvskill{Programmation}{Python, Java, JavaScript, C, SQL, Bash}
    \cvskill{Bases de données}{MySQL, PostgreSQL, Snowflake, Oracle DB, MongoDB, Vector DB (Milvus DB)}
    \cvskill{Machine Learning \& IA}{Scikit-learn, TensorFlow, LangChain \& LangGraph, Hugging Face, NLP (spaCy, NLTK), Flower FL}
    \cvskill{Ingénierie des données}{PySpark, Apache Kafka, Airflow, DBT, Hadoop Ecosystem (Couldera)}
     \cvskill{Analyse \& Reporting}{Excel, Power BI, Pandas, NumPy, Matplotlib, Seaborn, Plotly}
    \cvskill{Cloud \& DevOps}{AWS, GCP, OCI, Docker, MLOps, FastAPI, Git, Linux}
\end{cvskills}

%------------------------------------------------------------------
\vspace{-0.2cm}
\cvsection{Certifications}
\vspace{0.1cm}
\begin{cvhonors}
  \cvhonor
    {Oracle Cloud Infrastructure 2025 Certified Data Science Professional}
    {Oracle University}
    {}
    {2025}
  \cvhonor
    {Oracle Cloud Infrastructure 2025 Certified AI Foundations Associate}
    {Oracle University}
    {}
    {2025}
\end{cvhonors}


%-------------------------------------------------------------------------------
\vspace{-0.2cm}
\cvsection{Informations Complémentaires}
\vspace{0.1cm}
\begin{cvskills}
    \cvskill{Langues}{Arabe (Natif), Français (Courant), Anglais (Courant)}
    \cvskill{Points Forts}{Résolution de problèmes, Pensée analytique, Travail d'équipe, Apprentissage rapide}
\end{cvskills}
\end{document}
